% 三视图点修: 仅横向并排(主/左/俯左至右)+图名下置居中基线对齐; 投影边线/尺寸/虚实/比例均沿用既有稿
\begin{tikzpicture}[line width=.6pt,font=\normalsize,>=Latex]
% 主视图
\begin{scope}[shift={(0.6,1.2)}]
\draw (0,0) -- (1,0) -- (0,2) -- cycle;
% 底部宽1
\draw[thin] (0,0) -- (0,-0.5);
\draw[thin] (1,0) -- (1,-0.5);
\draw[thin,<->] (0,-0.3) -- (1,-0.3);
\node[below] at (0.5,-0.3) {1};
% 左侧高2
\draw[thin] (0,0) -- (-0.5,0);
\draw[thin] (0,2) -- (-0.5,2);
\draw[thin,<->] (-0.3,0) -- (-0.3,2);
\node[left] at (-0.3,1) {2};
\end{scope}
\node at (1.1,0) {正视图};
% 左视图
\begin{scope}[shift={(3.4,1.2)}]
\draw (-1,0) -- (1,0) -- (0,2) -- cycle;
\draw[thin] (-1,0) -- (-1,-0.5);
\draw[thin] (0,0) -- (0,-0.5);
\draw[thin] (1,0) -- (1,-0.5);
\draw[thin,<->] (-1,-0.3) -- (0,-0.3);
\draw[thin,<->] (0,-0.3) -- (1,-0.3);
\node[below] at (-0.5,-0.3) {1};
\node[below] at (0.5,-0.3) {1};
\end{scope}
\node at (3.4,0) {侧视图（左）};
% 俯视图: 左侧直径竖直,右侧半圆
\begin{scope}[shift={(5.2,2.2)}]
\draw (0,-1) -- (0,1);
\draw (0,1) arc (90:-90:1);
\end{scope}
\node at (5.7,0) {俯视图};
\end{tikzpicture}
